\documentclass[11pt,a4paper]{article}

\def\courseTitle{Industrial Security (Bachelor)}
\def\sectionNumber{05}
\def\sectionTitle{Threat Landscape}
\def\sheetKind{Exercise Sheet}

% =====================================================================
% Shared preamble for exercise sheets and solution sheets.
%
% Define BEFORE \input{}-ing this file:
%   \def\courseTitle{...}
%   \def\sectionNumber{...}
%   \def\sectionTitle{...}
%   \def\sheetKind{Exercise Sheet}   or  Solution Sheet
%
% Optional:
%   \def\courseAuthor{...}
%
% The caller must issue \documentclass{...} (typically article) BEFORE
% loading this preamble.
% =====================================================================

\usepackage[utf8]{inputenc}
\usepackage[T1]{fontenc}
\usepackage[english]{babel}
\usepackage[a4paper, margin=2.2cm]{geometry}
\usepackage{graphicx}
\usepackage{listings}
\usepackage{xcolor}
\usepackage{tikz}
\usepackage{booktabs}
\usepackage{array}
\usepackage{enumitem}
\usepackage{titlesec}
\usepackage{fancyhdr}
\usepackage{hyperref}
\usepackage{amsmath}
\usepackage{amssymb}
\usepackage{tcolorbox}
\tcbuselibrary{skins, breakable}
\usepackage{needspace}

% Discourage solo lines split across pages.
\widowpenalty=10000
\clubpenalty=10000
\raggedbottom

% =====================================================================
% Shared TikZ styles for the Industrial Security lecture series.
% Loaded by every slide deck and exercise sheet that uses TikZ.
% =====================================================================

\usetikzlibrary{
  arrows.meta,
  shapes,
  shapes.geometric,
  shapes.symbols,
  shapes.multipart,
  positioning,
  fit,
  calc,
  backgrounds,
  decorations.pathmorphing,
  decorations.pathreplacing,
  decorations.text,
  patterns,
  matrix,
  shadows
}

% --- colour palette --------------------------------------------------
\definecolor{isOT}{HTML}{1F4E79}     % OT (deep blue)
\definecolor{isIT}{HTML}{6F2DA8}     % IT (purple)
\definecolor{isSafety}{HTML}{C00000} % safety red
\definecolor{isNet}{HTML}{2E7D32}    % network green
\definecolor{isAttack}{HTML}{B23A48} % attack red
\definecolor{isAccent}{HTML}{E08E0B} % amber
\definecolor{isMute}{HTML}{6B6B6B}   % grey
\definecolor{isLight}{HTML}{EDF2F8}  % light background
\definecolor{isPLC}{HTML}{2C5282}
\definecolor{isHMI}{HTML}{2F855A}
\definecolor{isSCADA}{HTML}{6B46C1}

% --- generic node styles --------------------------------------------
\tikzset{
  isnode/.style={
    draw, rounded corners=2pt, minimum height=8mm, minimum width=18mm,
    font=\small, align=center, thick
  },
  isbox/.style={isnode, fill=isLight},
  plc/.style={isnode, fill=isPLC!15, draw=isPLC, thick},
  hmi/.style={isnode, fill=isHMI!15, draw=isHMI, thick},
  scada/.style={isnode, fill=isSCADA!15, draw=isSCADA, thick},
  sensor/.style={isnode, fill=isNet!10, draw=isNet, minimum height=6mm, minimum width=14mm, font=\scriptsize},
  actuator/.style={isnode, fill=isAccent!15, draw=isAccent, minimum height=6mm, minimum width=14mm, font=\scriptsize},
  firewall/.style={isnode, fill=isAttack!15, draw=isAttack, thick},
  server/.style={isnode, fill=isMute!15, draw=isMute!80, thick},
  switch/.style={isnode, fill=isLight, draw=black!60, minimum height=6mm, minimum width=14mm, font=\scriptsize},
  workstation/.style={isnode, fill=white, draw=isOT, thick},
  attacker/.style={isnode, fill=isAttack!20, draw=isAttack, thick, font=\small\bfseries},
  inetnode/.style={draw, ellipse, minimum width=24mm, minimum height=10mm, font=\scriptsize, fill=isLight, thick},
  process/.style={isnode, fill=isOT!15, draw=isOT, thick, minimum width=24mm},
  zone/.style={
    draw, rounded corners=4pt, thick, dashed, inner sep=4mm
  },
  conduit/.style={-{Stealth[length=2.5mm]}, thick, isNet!80!black},
  attack/.style={-{Stealth[length=2.5mm]}, thick, isAttack, dashed},
  flow/.style={-{Stealth[length=2.5mm]}, thick, isOT!70!black},
  netline/.style={thick, black!60},
  axisline/.style={-{Stealth[length=2mm]}, thick, black!70},
  tag/.style={font=\scriptsize\itshape, fill=white, inner sep=1pt},
  level/.style={
    draw, fill=isLight, minimum width=0.85\linewidth, minimum height=8mm,
    font=\small, anchor=west
  },
  brace/.style={decorate, decoration={brace, amplitude=4pt, raise=2pt}},
  legendbox/.style={draw, rounded corners=2pt, fill=isLight, inner sep=2mm, font=\scriptsize, align=left}
}

% --- handy macros ----------------------------------------------------
% \plcicon, \hmiicon etc as simple labelled boxes; useful inline.
\newcommand{\plcicon}[1][PLC]{\tikz[baseline=-0.6ex]{\node[plc, minimum width=10mm, minimum height=5mm, font=\scriptsize, inner sep=1pt]{#1};}}
\newcommand{\hmiicon}[1][HMI]{\tikz[baseline=-0.6ex]{\node[hmi, minimum width=10mm, minimum height=5mm, font=\scriptsize, inner sep=1pt]{#1};}}
\newcommand{\fwicon}[1][FW]{\tikz[baseline=-0.6ex]{\node[firewall, minimum width=10mm, minimum height=5mm, font=\scriptsize, inner sep=1pt]{#1};}}


\providecommand{\courseAuthor}{Industrial Security Lecture Series}
\providecommand{\courseInstitute}{Department of Computer Science}
\providecommand{\companyLogo}{logo}

% --- Corporate branding overrides ------------------------------------
% Picks up logo / author / brand colours from the repo-root branding.tex
% when present; falls back to the defaults above otherwise.
\IfFileExists{../../../branding.tex}{% =====================================================================
% Corporate / institutional branding -- single file to customise the
% look of the entire course (slides, notes, exercises, labs).
%
% HOW TO REBRAND IN UNDER A MINUTE:
%   1. Replace `branding/logo.pdf` with your own logo
%      (PDF preferred; PNG and JPG also work -- name it `logo.<ext>`).
%   2. (Optional) change the two colours below to your brand palette.
%   3. (Optional) override the author/institute lines below.
%   4. Re-run `make` at the repo root. That's it.
%
% Everything in this file has sensible defaults; you can delete a line
% to fall back to the built-in defaults, or comment it out.
%
% This file is loaded automatically by the slides, exercise, and lab
% preambles -- you never need to \input it manually.
% =====================================================================

% --- Logo --------------------------------------------------------------
% Path is resolved from each leaf-build directory; the default points at
% the shared `branding/` folder at the repo root. If you put your logo
% somewhere else, set the full or relative path here (without extension):
\renewcommand{\companyLogo}{../../../branding/logo}

% --- Author / institute (appears on the title page footer) ------------
\renewcommand{\courseAuthor}{}
\renewcommand{\courseInstitute}{}

% --- Brand colours ----------------------------------------------------
% slideAccent      = primary brand colour (title bands, accent rule)
% slideSecondary   = secondary brand colour (section badge, highlight)
% Both use HTML hex codes. Keep enough contrast against white text.
\definecolor{slideAccent}{HTML}{00205B}      % deep navy
\definecolor{slideSecondary}{HTML}{00A0DC}   % light blue accent
}{}

% --- listings -------------------------------------------------------
\definecolor{codebg}{rgb}{0.96,0.96,0.96}
\lstset{
  backgroundcolor=\color{codebg},
  basicstyle=\ttfamily\footnotesize,
  breaklines=true,
  frame=single,
  numbers=left,
  numberstyle=\tiny\color{black!50},
  showstringspaces=false,
  tabsize=2
}

% --- header / footer ------------------------------------------------
\pagestyle{fancy}
\fancyhf{}
\renewcommand{\headrulewidth}{0.4pt}
\renewcommand{\footrulewidth}{0pt}
\lhead{\small \courseTitle}
\rhead{\small Section \sectionNumber{} -- \sheetKind}
\cfoot{\thepage}

% --- task / solution boxes -----------------------------------------
% `before={...}` guards every task/solution box with \needspace so the
% box doesn't start with only one or two lines of breathing room at the
% bottom of a page. If less than ~9 lines remain, LaTeX bumps the box
% to the next page; otherwise it starts here and breaks normally if it
% runs long.
% Boxes are `breakable` so very long solutions don't blow the page,
% but `before={\needspace{14\baselineskip}}` pushes them to a fresh
% page whenever fewer than ~14 lines remain. That covers the vast
% majority of single-question splits in practice.
\newtcolorbox{taskbox}[1]{
  enhanced, breakable,
  colback=isLight, colframe=isOT,
  fonttitle=\bfseries\color{white},
  coltitle=white,
  colbacktitle=isOT,
  title={#1},
  arc=2pt, boxrule=0.6pt,
  left=2mm, right=2mm, top=2mm, bottom=2mm,
  before={\par\vspace{2mm}\needspace{14\baselineskip}\par},
  after={\par\vspace{2mm}}
}

\newtcolorbox{solutionbox}{
  enhanced, breakable,
  colback=isNet!5, colframe=isNet,
  fonttitle=\bfseries\color{white},
  coltitle=white,
  colbacktitle=isNet,
  title={Solution},
  arc=2pt, boxrule=0.6pt,
  left=2mm, right=2mm, top=2mm, bottom=2mm,
  before={\par\vspace{2mm}\needspace{14\baselineskip}\par},
  after={\par\vspace{2mm}}
}

\newtcolorbox{hintbox}{
  enhanced, breakable,
  colback=isAccent!10, colframe=isAccent,
  fonttitle=\bfseries\color{white}, coltitle=white, colbacktitle=isAccent,
  title={Hint},
  arc=2pt, boxrule=0.6pt
}

\newtcolorbox{warnbox}{
  enhanced, breakable,
  colback=isAttack!8, colframe=isAttack,
  fonttitle=\bfseries\color{white}, coltitle=white, colbacktitle=isAttack,
  title={Caution},
  arc=2pt, boxrule=0.6pt
}

% --- section formatting --------------------------------------------
\titleformat{\section}{\Large\bfseries\color{isOT}}{\thesection}{0.5em}{}
\titleformat{\subsection}{\large\bfseries\color{isOT!80!black}}{\thesubsection}{0.5em}{}

% --- title block macro ----------------------------------------------
\newcommand{\sheetheader}{%
  \begin{center}
    {\Large\bfseries \courseTitle}\\[0.3em]
    {\large \sheetKind{} -- Section \sectionNumber}\\[0.2em]
    {\large \sectionTitle}\\[0.2em]
    \rule{\linewidth}{0.4pt}
  \end{center}
  \vspace{0.5em}
}


\begin{document}
\sheetheader

\noindent
\textbf{Goal.} These exercises ask you to reason about real ICS incidents -- Stuxnet, Ukraine 2015/2016,
TRITON, NotPetya, Colonial, PIPEDREAM, CyberAv3ngers and Volt Typhoon -- using the case-study
material from the slides and the cited CISA / vendor / academic sources. Most questions are
short scenarios; the last one is an open mini-essay. Answers should be concise (3--10 sentences
each unless stated otherwise) and grounded in the published sources, not speculation.

\begin{taskbox}{Exercise 5.1 -- Stuxnet on the ATT\&CK for ICS Matrix}
Map the Stuxnet kill-chain to the MITRE ATT\&CK for ICS matrix. For each of the following six
phases of the attack, name (a)~the ATT\&CK \emph{tactic column} it sits in, and (b)~one concrete
\emph{technique ID} (T0xxx) you would assign:
\begin{enumerate}\setlength{\itemsep}{0pt}
  \item Crossing the air gap on a contractor's USB stick.
  \item Spreading inside the engineering LAN using a Windows print-spooler zero-day.
  \item Loading a kernel driver signed with a stolen Realtek certificate.
  \item Injecting malicious ladder logic into the S7-315 PLC via the Step7 DLL.
  \item Driving the Vacon / Fararo Paya VFDs outside their operating envelope.
  \item Replaying 21~seconds of recorded normal telemetry to the HMI during the attack.
\end{enumerate}
Present your answer as a three-column table (\emph{Phase, Tactic, Technique~ID}). Cite the
ATT\&CK for ICS matrix as your reference.
\end{taskbox}

\begin{taskbox}{Exercise 5.2 -- Ukraine 2015 vs.\ Ukraine 2016}
Both the December 2015 BlackEnergy attack and the December 2016 Industroyer / CrashOverride
attack blacked out parts of Ukraine, but they were qualitatively different operations.
Compare them in a short table covering these four axes:
\begin{enumerate}\setlength{\itemsep}{0pt}
  \item \textbf{Target tier} -- distribution oblenergos vs.\ transmission substation.
  \item \textbf{Human-in-the-loop} -- did a human click breakers, or did malware?
  \item \textbf{Protocol awareness} -- generic remote control vs.\ purpose-built modules
        (IEC 60870-5-101/104, IEC 61850 MMS, OPC DA).
  \item \textbf{Recovery effort / time} -- outage hours and what restoration required.
\end{enumerate}
Conclude with one sentence on what the 2015~$\to$~2016 transition implies about attacker capability.
\end{taskbox}

\begin{taskbox}{Exercise 5.3 -- TRITON: a Partial Defender Win}
TRITON / TRISIS (Saudi petrochemical plant, August~2017) targeted a Schneider Triconex safety
instrumented system (SIS). In production, the SIS tripped the plant to a safe state and the
campaign was exposed.
\begin{enumerate}\setlength{\itemsep}{0pt}
  \item In 2--3 sentences, describe what the attacker presumably \emph{wanted} the
        SIS to do, and how this was meant to combine with a separate attack on the BPCS.
  \item Explain in 2--3 sentences why the actual outcome (safe-state trip + discovery)
        was a partial defender win, and why it was nonetheless deeply unsettling.
  \item Name one technical reason TRITON-class implants are hard to detect on
        Triconex controllers (where the malicious logic lived, what protocol it spoke).
\end{enumerate}
\end{taskbox}

\begin{taskbox}{Exercise 5.4 -- Volt Typhoon and Living off the Land}
CISA advisory \emph{AA24-038A} (7~Feb~2024) attributes a long-running campaign against US critical
infrastructure to Volt Typhoon, who rely almost exclusively on \emph{living-off-the-land binaries}
(LOLBins) instead of custom malware.
\begin{enumerate}\setlength{\itemsep}{0pt}
  \item List \textbf{four} LOLBins or built-in Windows tools that the slides and AA24-038A
        associate with this actor (e.g.\ \texttt{wmic}, \texttt{ntdsutil}, PowerShell, \ldots).
  \item For each, explain in one sentence what an admin would normally use it for.
  \item Explain in 3--5~sentences why a defender's signature-based endpoint product (AV / EDR
        signatures keyed on file hashes or known IOCs) is structurally unable to flag this kind
        of activity, and what class of detection \emph{can}.
\end{enumerate}
\end{taskbox}

\begin{taskbox}{Exercise 5.5 -- PIPEDREAM: Caught Before Impact}
PIPEDREAM / INCONTROLLER / CHERNOVITE (CISA \emph{AA22-103A}, 13~Apr~2022) is the first modular
ICS attack framework that was discovered \emph{before} it caused any victim impact, thanks to
early reporting by a partner.
\begin{enumerate}\setlength{\itemsep}{0pt}
  \item Why does ``caught before deployment'' matter \emph{geopolitically}, compared to a
        post-incident discovery like Stuxnet or Industroyer? (2--3 sentences.)
  \item What does PIPEDREAM's design -- pre-built modules for Schneider M340/M580, Omron NX/NJ,
        OPC~UA and CODESYS, with scanning, parameter changes, logic upload and denial of
        control -- tell you about the maturity of state-sponsored ICS tooling in 2022? (2--3
        sentences.)
  \item Name one defensive action a Schneider or Omron asset owner should have taken
        \emph{immediately} on reading AA22-103A.
\end{enumerate}
\end{taskbox}

\begin{taskbox}{Exercise 5.6 -- Colonial Pipeline: Decoupling IT and OT}
The Colonial Pipeline incident (7~May~2021) is widely cited as an ``IT incident that took OT
down''. The DarkSide affiliate compromised the corporate IT environment; OT control was never
directly attacked, yet Colonial halted the 8\,850\,km pipeline for five days.
\begin{enumerate}\setlength{\itemsep}{0pt}
  \item In your own words, describe the dependency chain that turned an IT ransomware event
        into a fuel-supply crisis on the US east coast. Be specific about \emph{which} IT
        function the operator could not do without.
  \item Propose \textbf{two} architectural or organisational changes that would have allowed
        OT to keep flowing product even while corporate IT was encrypted. For each, name one
        cost or trade-off (capex, complexity, regulatory, $\ldots$) that explains why it was
        not in place.
\end{enumerate}
\end{taskbox}

\begin{taskbox}{Exercise 5.7 -- Reading a CISA Advisory: AA23-335A}
Read CISA advisory \emph{AA23-335A} (``IRGC-Affiliated Cyber Actors Exploiting PLCs in Multiple
Sectors, Including U.S.\ Water and Wastewater Systems Facilities'', 1~Dec~2023). Answer:
\begin{enumerate}\setlength{\itemsep}{0pt}
  \item Which specific PLC product family did CyberAv3ngers target, and on which TCP port does
        its engineering protocol listen?
  \item Which default password did the attackers use, and what was the highest-profile
        US~victim cited in the advisory?
  \item What is CISA's \emph{first} listed mitigation recommendation? (Quote or paraphrase the
        action verb.)
  \item In one sentence: which one of the eight initial-access vectors on the slides was the
        decisive one for this campaign?
\end{enumerate}
\end{taskbox}

\begin{taskbox}{Exercise 5.8 -- Pyramid of Pain}
The ``Pyramid of Pain'' (Bianco, 2013) ranks indicators of compromise by how painful they are
for the attacker to change. Rank the following five IOC types from \textbf{easiest} (attacker
swaps in minutes) to \textbf{hardest} (forces the attacker to rebuild their tradecraft), and
justify each placement in one sentence:
\begin{itemize}\setlength{\itemsep}{0pt}
  \item MD5 hash of a payload binary.
  \item C2 IP address.
  \item Registry key used for autorun persistence.
  \item C2 domain name.
  \item MITRE ATT\&CK technique (e.g.\ T0855 ``Unauthorized Command Message'').
\end{itemize}
Conclude with one sentence on what this implies for a SOC that detects \emph{only} on hashes.
\end{taskbox}

\begin{taskbox}{Exercise 5.9 -- Insider Threat: First 24 Hours}
A contract control-systems engineer at a mid-size chemical plant hands in their notice on a
Friday afternoon. On the same day, SIEM logs show that their account downloaded the full DCS
project file (HMI graphics, controller config, tag database, alarm setpoints) from the
engineering server to a personal laptop over the corporate VPN, outside their usual
working hours.

Walk through what the security team should do in the \textbf{first 24~hours}. Cover, in order:
\begin{enumerate}\setlength{\itemsep}{0pt}
  \item Containment actions on accounts, VPN sessions and physical access.
  \item Evidence preservation -- which artefacts must be captured before they age out?
  \item Who to notify inside the organisation (and roughly when), and which external parties
        (legal, HR, law enforcement, sectoral CERT) come into scope and on what trigger.
  \item One thing the team must \emph{not} do that would be tempting (think evidentiary
        integrity, employment law, or alerting the suspect).
\end{enumerate}
\end{taskbox}

\begin{taskbox}{Exercise 5.10 -- Open Research: A Decision You Would Reverse}
Choose \textbf{one} of the case-study incidents covered in this section -- Stuxnet, Ukraine
2015, Ukraine 2016 (Industroyer), TRITON, NotPetya, Colonial Pipeline, PIPEDREAM,
CyberAv3ngers, or Volt Typhoon -- and write a one-page (about 400--600~word) analysis of
\emph{one} decision the defenders (asset owner, integrator, vendor, or regulator) made before
or during the incident that, with hindsight, you would do differently.

Your write-up should: (i)~name the decision and the actor who made it, citing at least one
primary source (CISA advisory, ICS-CERT alert, vendor PSIRT, peer-reviewed paper, or named
journalism / book); (ii)~explain the operational or commercial reasoning that made the
decision defensible at the time; (iii)~describe the alternative you would have chosen and
the cost it would have carried; (iv)~assess whether your alternative would have changed the
outcome materially, or only at the margin.
\end{taskbox}

\end{document}
